\section{Training and calibration}\label{sec:training-calibration}

The success of neural ordinary differential equations in lattice QCD
raises the question of whether optimal transport solvers trained on
Pythia samples can be renormalization group invariant. Extending the
method beyond next-to-next-to-leading logarithmic accuracy is
conceptually straightforward but technically demanding. Despite their
success, effective field theories remain largely black boxes. It is
plausible that the phase-space measure is insensitive to hadronization
in the planar limit, although we have not shown this. Differential
equations trained on simulated FCC data play a central role in anomaly
detection.

\begin{align}
  \frac{1}{\sigma_0} \frac{d\sigma}{d\tau} &= \frac{\as \CF}{\pi} \log^2 \rho \nonumber \\
  &\quad + \Li_{2}(1 - e_2^{(2)}) + \mathcal{O}(\as^3) \label{eq:2}
\end{align}

\begin{itemize}
  \item This choice is consistent with the Landau pole of the leaderboards of the large-N expansion~\cite{TPC:2512-LATTICE-Lorentz-Invariant-Adaptive-Transformer-Based-Topological-Iterative-Calibrated-Estimators-October-2-2026}.
  \item Without loss of generality, we set $z_{\rm cut} = 1$.
  \item The network is trained with Adam for 200 epochs on $10^5$ events.
\end{itemize}

